\subsection{A pole-aware pool member: a partly negative result}\label{sec:rational}
The two Helmholtz problems above fail for a structural reason: the response has poles inside the parameter range, which Gaussian processes, kernel ridge and ridge polynomials cannot represent. We therefore tried a method contribution aimed at exactly that case, and tested it as a pool member of the tool. \texttt{PodRational} regresses the POD coefficients $c_k(\mu)$ by a rational function with a shared denominator, $c_k=P_k(\mu)/Q(\mu)$, with $P_k$ and $Q$ polynomials of low total degree. The denominator is fitted by Sanathanan--Koerner iterations~\cite{sanathanan1963} on the leading modes (variable projection, ridge), the numerators of all modes by weighted ridge, and the degrees and the ridge are chosen on a 25\% hold-out of the training runs from a fixed grid.
\paragraph{Relation to earlier work} Rational approximation of frequency-response maps is established: the p-AAA algorithm builds multivariate barycentric rational surrogates of parametric dynamical systems~\cite{carracedo2023}, least-squares Pad\'e approximation treats parametric and stochastic Helmholtz maps~\cite{bonizzoni2020}, and non-intrusive greedy rational model reduction uses rational surrogates in the frequency domain~\cite{nobile2021}. What is tested here is narrower and more of an engineering combination than a new theory: a data-driven shared-denominator regression of POD coefficients with hold-out-selected degrees, entered in a cross-validated pool next to Gaussian processes and networks. Whether that is more than an incremental combination is for the reader to judge; our literature search for the combination was limited.
\paragraph{Design and criteria} The member was developed on seen data only (a synthetic exact rational response, in which it recovers the truth to $3\cdot10^{-5}$ where a polynomial gives 0.3, and the seen Helmholtz and plate problems). Its test was then fixed in writing (addendum v3.5, freeze commit in Supplement Sec.~\ref{sec:S-prereg}) on four new, unseen frequency-response problems with damping (complex linear systems; real and imaginary parts as outputs; Supplement Table~\ref{tab:S-srproblems}): an acoustic cavity with impedance walls (square and L-shaped) and the frequency response of a damped cantilever with a tip mass. Their solvers pass checks against theory (the square cavity resonates at $k=\pi$, the cantilever at the Euler--Bernoulli value, dissipation is non-negative, the system is reciprocal, the mesh converges). Six criteria were set before any fit (\cref{tab:srcells} for the data).
\begin{table}[t]\centering\caption{Frequency-response problems: test error (geometric mean over three seeds) of the model selected by the tool with the old pool (seven classical members), with the pool that includes the rational member, and of the tuned network; and how often the rational member was chosen.}\label{tab:srcells}\small\resizebox{\textwidth}{!}{\begin{tabular}{lrrrrr}
\toprule
Problem & $N$ & Old pool & With rational & Tuned net & Rational chosen \\
\midrule
\texttt{acoustic\_cavity\_L1} & 256 & $5.2\times10^{-2}$ & $2.3\times10^{-2}$ & $2.8\times10^{-2}$ & 3 of 3 \\
\texttt{acoustic\_cavity\_L1} & 1024 & $1.4\times10^{-2}$ & $2.8\times10^{-3}$ & $1.5\times10^{-2}$ & 3 of 3 \\
\texttt{acoustic\_cavity\_L2} & 256 & $1.3\times10^{-1}$ & $1.3\times10^{-1}$ & $1.1\times10^{-1}$ & 0 of 3 \\
\texttt{acoustic\_cavity\_L2} & 1024 & $4.2\times10^{-2}$ & $2.2\times10^{-2}$ & $4.2\times10^{-2}$ & 3 of 3 \\
\texttt{beam\_frf\_L1} & 256 & $1.3\times10^{0}$ & $1.6\times10^{0}$ & $1.9\times10^{-1}$ & 1 of 3 \\
\texttt{beam\_frf\_L1} & 1024 & $8.0\times10^{-1}$ & $4.0\times10^{-1}$ & $1.4\times10^{-1}$ & 3 of 3 \\
\texttt{beam\_frf\_L2} & 256 & $2.7\times10^{0}$ & $2.3\times10^{0}$ & $2.2\times10^{0}$ & 1 of 3 \\
\texttt{beam\_frf\_L2} & 1024 & $1.8\times10^{0}$ & $1.4\times10^{0}$ & $8.7\times10^{-1}$ & 2 of 3 \\
\bottomrule
\end{tabular}
}\end{table}
\paragraph{Outcome} \SrMetCount{} of the six criteria were met. The gain criterion was \emph{not} met: the error was at least halved in only \SrOneK{} of \SrOneN{} cells, where we had required six (R1 \SrOneVerdict). The gains are real where the structure fits: on the square cavity the selected error fell from \SrATwofiftysixOld{} to \SrATwofiftysixNew{} at 256 runs and from \SrATentwentyfourOld{} to \SrATentwentyfourNew{} at 1024, and on the cantilever with one resonance in range from \SrBTentwentyfourOld{} to \SrBTentwentyfourNew{} at 1024; the error improved by more than 20\% in \SrFBetterFifth{} cells, all with at least 256 runs, none at 64. The member never lowered the accuracy by more than a quarter (R2 \SrTwoVerdict: \SrTwoWorse{} cell worse by more than 10\%, a wrong choice by cross-validation at 256 runs on the lightly damped cantilever), and the selection stayed good (R3 \SrThreeVerdict: median regret \SrThreeMedian). It did not make the models usable: the number of cells with error at most 0.25 was \SrFiveNew{} of \SrFiveN{}, the same as before (R5 \SrFiveVerdict). Most importantly, the tuned network was better than the pool with the rational member by more than 25\% in \SrFourK{} of \SrFourN{} cells (R4 \SrFourVerdict), and on the cantilever with one resonance in range it was the only usable model (\cref{tab:srcells}). A likely reason, which we did not test, is that a rational denominator of total degree four in four inputs cannot represent a resonance whose position, width and height depend on frequency, damping, mass and load position, whereas a network can approximate that map. On the \SrSixN{} earlier problems the member did no harm (median ratio \SrSixMedian, \SrSixWorse{} problems worse by more than 10\%; R6 \SrSixVerdict); it was selected on one of them, \texttt{helmholtz\_L1}, where it cut the error from \SrHHelmOld{} to \SrHHelmNew{}. Its cost is small (about \SrRatShare{} of the tool's cross-validation time; \SrRatSecTentwentyfour{} seconds at 1024 runs).
\paragraph{What we conclude} The rational member does what it was designed for in a narrow case: lightly damped systems with a single resonance in the parameter range and at least a few hundred runs, where it removes most of the error that the classical members cannot. It misses the criterion we set for a general improvement, it loses to a tuned network where the dependence is richer, and our predictions of R1 and R5 were wrong. By the rule in the addendum it is not in the default pool of tool version 0.4.1; it is an opt-in option (\texttt{--with-rational}). The practical advice that follows from this stage and from \cref{sec:confirm} is consistent: when the classical pool reports that no model is usable, add the network.
